\documentclass[a4paper,12pt]{article}

\usepackage[T2A]{fontenc}
\usepackage[utf8]{inputenc}
\usepackage[english,russian]{babel}
\usepackage{geometry}
\usepackage{setspace}
\usepackage{amsmath}
\usepackage{amssymb}
\usepackage{amsthm}
\usepackage{array}
\usepackage{booktabs}
\usepackage{enumitem}
\usepackage{hyperref}

\geometry{left=2.5cm,right=1.8cm,top=2cm,bottom=2cm}
\setlength{\parindent}{1.25cm}
\setstretch{1.35}
\hypersetup{colorlinks=true,linkcolor=black,urlcolor=blue,citecolor=black}

\newtheorem{definition}{Определение}[section]
\newtheorem{lemma}{Лемма}[section]
\newtheorem{proposition}{Предложение}[section]
\newtheorem{corollary}{Следствие}[section]

\newcommand{\dist}{d}
\newcommand{\obj}{F}
\newcommand{\clients}{V \setminus \{0\}}

\begin{document}

\begin{center}
{\Large Приложение. Математическое обоснование реализованных улучшений для задачи mTSP}

\vspace{0.5em}
{\large для проекта \texttt{mtsp-routing-optimization}}
\end{center}

\vspace{1em}

\section{Назначение приложения}

Настоящее приложение фиксирует математические основания тех эвристических методов, которые
реально реализованы в проекте \texttt{mtsp-routing-optimization}. Речь идет не об абстрактном
обзоре по задаче множественного коммивояжера вообще, а о формальном описании и
обосновании алгоритмов, вынесенных в программные модули
\path{src/mtsp_rand_nn_solver.cpp},
\path{src/mtsp_greedy_2opt_solver.cpp},
\path{src/mtsp_grasp_solver.cpp} и
\path{src/mtsp_lkh_wrapper_v2.cpp}.

Главная цель приложения состоит в том, чтобы показать, почему реализованные преобразования
сохраняют допустимость решения, почему процедуры локального улучшения завершаются за
конечное число шагов и за счет каких именно механизмов более сильные эвристики получают
преимущество над простыми baseline-подходами.

\section{Постановка задачи}

В проекте рассматривается замкнутый single-depot вариант задачи множественного
коммивояжера с фиксированным числом агентов и целевой функцией \texttt{MINSUM}. Пусть
\[
V = \{0,1,\dots,n-1\},
\]
где вершина \(0\) является депо, а множество клиентов равно \(\clients\).
Каждой вершине \(i \in V\) поставлена в соответствие точка
\((x_i, y_i) \in \mathbb{R}^2\), а расстояние между вершинами определяется евклидовой метрикой
\[
\dist(i,j) = \sqrt{(x_i - x_j)^2 + (y_i - y_j)^2}.
\]

Пусть \(m \ge 1\) --- число коммивояжеров. Решением задачи является набор маршрутов
\[
\mathcal{R} = (R_1,\dots,R_m),
\]
где каждый маршрут имеет вид
\[
R_k = (0, v_{k,1}, v_{k,2}, \dots, v_{k,q_k}, 0).
\]

\begin{definition}
Набор маршрутов \(\mathcal{R}\) называется \emph{допустимым}, если выполняются два условия:
\begin{enumerate}[label=\arabic*.]
  \item каждый маршрут начинается и заканчивается в депо;
  \item каждая вершина из \(\clients\) входит ровно в один маршрут и встречается в нем ровно один раз.
\end{enumerate}
\end{definition}

Длина отдельного маршрута задается формулой
\[
L(R_k) = \sum_{t=0}^{q_k} \dist(v_{k,t}, v_{k,t+1}),
\]
где по соглашению \(v_{k,0} = v_{k,q_k+1} = 0\).
Целевая функция задачи имеет вид
\[
\obj(\mathcal{R}) = \sum_{k=1}^{m} L(R_k).
\]
Требуется найти допустимое решение \(\mathcal{R}\), минимизирующее \(\obj(\mathcal{R})\).

Такая постановка совпадает с классическим вариантом mTSP с общим депо и критерием
минимизации суммарной длины маршрутов \cite{bektas2006}.

\section{Алгоритм \texttt{rand+nn}}

\subsection{Схема работы}

Алгоритм \texttt{rand+nn} состоит из двух шагов:
\begin{enumerate}[label=\arabic*.]
  \item множество клиентов случайно перемешивается и затем распределяется между агентами по
  схеме round-robin;
  \item для каждого полученного подмножества клиентов строится маршрут ближайшего соседа,
  который стартует из депо и в конце возвращается в депо.
\end{enumerate}

В программной реализации этому соответствует случайная перестановка массива вершин,
разбиение по индексам \(t \bmod m\) и вызов процедуры \texttt{BuildNearestOrder}.

\subsection{Корректность}

\begin{proposition}
Алгоритм \texttt{rand+nn} всегда строит допустимое решение задачи mTSP.
\end{proposition}

\begin{proof}
После случайного перемешивания каждый клиент из \(\clients\) присутствует в массиве ровно
один раз. Разбиение по схеме round-robin помещает каждый клиентский индекс ровно в одно из
\(m\) подмножеств, то есть получаем попарно непересекающееся покрытие множества
\(\clients\).

Далее для каждого подмножества \(S_k \subseteq \clients\) процедура ближайшего соседа
последовательно выбирает очередную непосещенную вершину из \(S_k\), добавляет ее в маршрут
и удаляет из пула доступных вершин. Следовательно, каждая вершина из \(S_k\) включается в
маршрут ровно один раз. В начало и конец маршрута явно добавляется депо \(0\).

Таким образом, каждый маршрут имеет вид
\((0,\dots,0)\), а объединение внутренних вершин по всем маршрутам совпадает с
\(\clients\) без повторений. Это и означает допустимость решения.
\end{proof}

\subsection{Оценка сложности}

Перемешивание клиентов требует \(O(n)\) времени. Построение маршрута ближайшего соседа
на подмножестве размера \(q_k\) требует \(O(q_k^2)\), поскольку на каждом шаге просматривается
весь оставшийся пул вершин. Поэтому суммарная сложность алгоритма оценивается как
\[
O\!\left(\sum_{k=1}^{m} q_k^2 \right) \subseteq O(n^2).
\]

Алгоритм дает корректное, но грубое начальное решение: распределение клиентов между
маршрутами не использует информацию о геометрии задачи, поэтому основная потеря качества
происходит еще до этапа упорядочения внутри каждого маршрута.

\section{Алгоритм \texttt{2opt+greed}}

\subsection{Схема работы}

Алгоритм \texttt{2opt+greed} также состоит из двух частей:
\begin{enumerate}[label=\arabic*.]
  \item в фазе построения каждый агент по очереди получает ближайший непосещенный город
  относительно своей текущей позиции;
  \item затем для каждого маршрута выполняется внутримаршрутный локальный поиск 2-opt
  до достижения локального минимума в 2-opt-окрестности.
\end{enumerate}

\subsection{Сохранение допустимости при 2-opt}

\begin{lemma}
Пусть маршрут \(R = (0,v_1,\dots,v_q,0)\) допустим. Если на нем выполнен 2-opt-обмен,
то полученный маршрут содержит те же вершины, начинается и заканчивается в депо, а значит,
остается допустимым.
\end{lemma}

\begin{proof}
2-opt удаляет две непересекающиеся дуги маршрута и заменяет их двумя новыми дугами,
соединяющими те же четыре граничные вершины в другом порядке \cite{croes1958}.
Эквивалентно это можно описать как разворот внутреннего сегмента между двумя позициями
маршрута. При таком развороте множество внутренних вершин не меняется, повторов и
потерь вершин не возникает, а крайние вершины маршрута остаются равными \(0\).
Следовательно, допустимость сохраняется.
\end{proof}

\subsection{Монотонность локального улучшения}

\begin{proposition}
Если 2-opt-обмен принимается только в случае строгого уменьшения длины маршрута, то
каждый принятый шаг строго уменьшает значение целевой функции \(\obj\).
\end{proposition}

\begin{proof}
В алгоритме \texttt{2opt+greed} межмаршрутное распределение после фазы построения не
изменяется. Следовательно, при локальном поиске меняется только длина одного маршрута
\(R_k\), а длины остальных маршрутов остаются постоянными. Если найден обмен,
уменьшающий \(L(R_k)\), то
\[
\obj(\mathcal{R}_{new}) =
\obj(\mathcal{R}_{old}) - \bigl(L(R_k^{old}) - L(R_k^{new})\bigr)
< \obj(\mathcal{R}_{old}).
\]
\end{proof}

\begin{corollary}
Алгоритм \texttt{2opt+greed} завершается за конечное число шагов и возвращает решение,
локально оптимальное относительно 2-opt-перестановок внутри каждого маршрута.
\end{corollary}

\begin{proof}
Множество допустимых решений на конечном наборе вершин конечно. Каждый принятый 2-opt
шаг строго уменьшает \(\obj\). Бесконечная строго убывающая последовательность на конечном
множестве невозможна, значит, процесс завершается. В финальной точке больше нет
улучшающего 2-opt-хода, следовательно, достигнут локальный минимум.
\end{proof}

\subsection{Оценка сложности}

Фаза жадного распределения выполняет \(n-1\) назначений, и каждое назначение в худшем
случае просматривает все еще не посещенные вершины. Это дает \(O(n^2)\).

Один полный проход 2-opt по маршруту длины \(q_k+2\) содержит \(O(q_k^2)\) проверок.
Если до остановки требуется \(I_k\) проходов, то суммарная стоимость локального улучшения
оценивается как
\[
O\!\left(\sum_{k=1}^{m} I_k q_k^2 \right).
\]

\section{Алгоритм \texttt{grasp}}

\subsection{Схема работы}

Алгоритм \texttt{grasp} реализует стандартную для GRASP двухфазную структуру
\cite{feo1995}:
\begin{enumerate}[label=\arabic*.]
  \item на каждой итерации строится жадно-рандомизированное начальное решение;
  \item полученное решение улучшается локальным поиском;
  \item из всех итераций сохраняется лучшее решение.
\end{enumerate}

В фазе построения для каждого агента формируется список кандидатов, отсортированных по
расстоянию до текущей вершины. Далее выбирается не обязательно лучший кандидат, а случайный
элемент из ограниченного списка ближайших соседей RCL.

\subsection{Допустимость и конечность}

\begin{proposition}
Каждая отдельная итерация \texttt{grasp} порождает допустимое решение.
\end{proposition}

\begin{proof}
В фазе построения используется массив \texttt{visited}. Как только клиент включается в один из
маршрутов, он помечается как посещенный и больше не участвует в выборе. Значит, каждый
клиент может быть назначен не более одного раза. Цикл продолжается до тех пор, пока число
неназначенных клиентов не станет нулевым, поэтому каждый клиент назначается ровно один раз.
После этого к каждому маршруту добавляется возврат в депо, а дальнейшие этапы локального
поиска используют либо 2-opt внутри маршрута, либо попарные обмены вершинами между
маршрутами, то есть сохраняют множество клиентов и не нарушают структуру
\((0,\dots,0)\).
\end{proof}

\begin{proposition}
Локальное улучшение внутри одной итерации \texttt{grasp} завершается за конечное число шагов.
\end{proposition}

\begin{proof}
После конструктивной фазы выполняются два типа улучшающих преобразований:
\begin{enumerate}[label=\arabic*.]
  \item 2-opt внутри маршрута;
  \item межмаршрутный обмен двух клиентов, принимаемый только при строгом уменьшении
  \(\obj\).
\end{enumerate}
Оба вида шагов сохраняют допустимость решения. Каждый принятый шаг строго уменьшает
целевую функцию, а число допустимых решений конечно. Следовательно, бесконечный процесс
невозможен.
\end{proof}

\subsection{Оценка сложности}

Пусть \(T\) --- число итераций GRASP, а размер RCL ограничен константой \(r\). На одной
итерации построение требует не менее \(O(n^2)\) времени, а при явной сортировке списков
кандидатов можно использовать консервативную оценку \(O(n^2 \log n)\).

Дальнейший локальный поиск включает:
\begin{enumerate}[label=\arabic*.]
  \item внутримаршрутный 2-opt со стоимостью
  \(O\!\left(\sum_{k=1}^{m} I_k q_k^2\right)\);
  \item межмаршрутные обмены, число которых зависит от конкретного экземпляра и
  траектории улучшения.
\end{enumerate}

Следовательно, полная трудоемкость \texttt{grasp} за все итерации оценивается сверху как
\[
O\!\left(T \cdot \left(n^2 \log n + \sum_{k=1}^{m} I_k q_k^2 + W\right)\right),
\]
где \(W\) обозначает стоимость межмаршрутных улучшений на одной итерации.

\section{Алгоритм \texttt{lkh-wrapper-v2}}

\subsection{Общая идея}

Алгоритм \texttt{lkh-wrapper-v2} не является прямым вызовом внешнего пакета LKH-3.
Он представляет собой собственную эвристику проекта, вдохновленную идеями
Lin--Kernighan и Helsgaun \cite{lin1973,helsgaun2000,russell1977}. По сравнению с
baseline-методами в нем реализованы четыре ключевых усиления:
\begin{enumerate}[label=\arabic*.]
  \item предварительное построение \emph{candidate sets} для каждой вершины;
  \item lookahead-оценка на этапе построения начального решения;
  \item candidate-guided 2-opt с последующим полным cleanup-проходом;
  \item итерационный локальный поиск с крупным возмущением и межмаршрутными обменами.
\end{enumerate}

\subsection{Lookahead-построение}

Для каждой вершины \(u \in V\) заранее выбирается множество кандидатов
\(\Gamma(u)\) из \(k\) ближайших соседей. Затем на этапе построения начального решения
для возможного перехода \(u \to v\) вычисляется оценка
\[
S(u,v) = \dist(u,v) + \lambda \,\phi(v) + \mu \,\dist(v,0),
\]
где \(\phi(v)\) --- оценка следующего выгодного продолжения из вершины \(v\),
\(\lambda > 0\) --- вес lookahead-компоненты, а \(\mu > 0\) --- вес штрафа за удаление от депо.

Такой критерий не минимизирует целевую функцию точно, но уменьшает вероятность ранних
неудачных назначений, при которых ближайший шаг выглядит выгодным локально, а дальнейшее
продолжение маршрута резко ухудшается. Именно это и является математическим смыслом
lookahead: выбор определяется не только текущим ребром, но и приближенной оценкой
следующего шага.

\begin{proposition}
Lookahead-конструкция в \texttt{lkh-wrapper-v2} сохраняет допустимость решения.
\end{proposition}

\begin{proof}
Как и в предыдущих алгоритмах, каждый клиент может быть добавлен в маршрут только пока
он не отмечен в массиве \texttt{visited}. После включения клиента в один из маршрутов эта
отметка становится истинной и клиент больше не рассматривается. Значит, каждая вершина из
\(\clients\) назначается не более одного раза. Цикл построения останавливается только после
обнуления счетчика оставшихся клиентов, то есть каждая вершина назначается ровно один раз.
Старт и завершение каждого маршрута в депо задаются явно.
\end{proof}

\subsection{Candidate-guided 2-opt}

Внутримаршрутный локальный поиск в \texttt{lkh-wrapper-v2} использует заранее построенные
candidate sets. Для текущего ребра \((t_1,t_2)\) рассматриваются только такие вершины
\(t_3 \in \Gamma(t_1)\), которые реально присутствуют в маршруте и могут образовать допустимый
2-opt-обмен. Это сокращает число проверяемых вариантов с полного перебора до перебора по
ограниченному списку кандидатов, после чего выполняется обычный cleanup-проход полного
2-opt.

\begin{lemma}
Candidate-guided 2-opt в \texttt{lkh-wrapper-v2} сохраняет допустимость маршрута.
\end{lemma}

\begin{proof}
Используемое преобразование остается тем же самым 2-opt-обменом: меняется лишь правило
выбора пар ребер, а не геометрия самой перестройки. Согласно лемме 4.1, любая корректно
выполненная 2-opt-замена сохраняет множество вершин маршрута и концевые вершины, равные
депо. Следовательно, ограничение поиска candidate sets не влияет на корректность, а влияет
только на ширину рассматриваемой окрестности.
\end{proof}

\begin{proposition}
Если candidate-guided 2-opt принимает только шаги с положительным выигрышем, то после
завершения процедуры маршрут является локально оптимальным относительно реализованной
ограниченной окрестности.
\end{proposition}

\begin{proof}
Каждый принятый шаг строго уменьшает длину текущего маршрута. Число перестановок конечного
маршрута конечно. Поэтому процедура завершается. В момент остановки не существует
кандидатного 2-opt-хода с положительным выигрышем, а после cleanup-прохода не существует
и улучшающего полного 2-opt-хода. Значит, достигнут локальный минимум в реализованной
окрестности.
\end{proof}

\subsection{Итерационный локальный поиск}

После локальной оптимизации маршрут подвергается крупному возмущению типа
double-bridge, а затем снова улучшается локальным поиском. Если новый вариант лучше,
он становится текущим лучшим. Это дает стандартную для iterated local search схему:
локальный минимум разрушается достаточно сильным преобразованием, после чего поиск
восстанавливает хорошую структуру решения уже в другом бассейне притяжения.

\begin{proposition}
Элитистское правило отбора в \texttt{lkh-wrapper-v2} делает последовательность лучших
найденных значений монотонно невозрастающей.
\end{proposition}

\begin{proof}
Пусть \(R^{(t)}\) --- лучшее решение после \(t\)-го раунда локального поиска. Если новое
кандидатное решение хуже, алгоритм оставляет \(R^{(t)}\) без изменения. Если кандидат лучше,
то оно заменяет текущее лучшее решение. Следовательно,
\[
\obj(R^{(t+1)}) \le \obj(R^{(t)})
\]
для всех \(t\), то есть последовательность лучших значений не возрастает.
\end{proof}

\subsection{Межмаршрутные обмены}

Финальный этап \texttt{lkh-wrapper-v2} перебирает пары маршрутов и пытается обменивать
внутренние вершины. Шаг принимается только тогда, когда суммарная длина набора маршрутов
строго уменьшается. После принятия обмена оба затронутых маршрута дополнительно
оптимизируются локальным поиском.

\begin{proposition}
Принятый межмаршрутный обмен в \texttt{lkh-wrapper-v2} сохраняет допустимость решения.
\end{proposition}

\begin{proof}
Обмен производится между двумя внутренними позициями разных маршрутов. Ни одна вершина
не создается и не удаляется: две вершины просто меняются местами. Поскольку до обмена каждая
вершина встречалась ровно один раз, то и после обмена это свойство сохраняется. Начало и
конец каждого маршрута, равные \(0\), не изменяются. Следовательно, допустимость решения
сохраняется.
\end{proof}

\begin{corollary}
Алгоритм \texttt{lkh-wrapper-v2} завершается за конечное число шагов и возвращает допустимое
решение, локально оптимальное относительно совокупности реализованных внутримаршрутных и
межмаршрутных преобразований.
\end{corollary}

\begin{proof}
Все принимаемые изменения в алгоритме либо уменьшают длину одного маршрута, либо
уменьшают суммарную длину всей системы маршрутов. При этом допустимость решения
сохраняется. Множество допустимых решений конечно, следовательно, бесконечная цепочка
строгих улучшений невозможна.
\end{proof}

\subsection{Оценка сложности}

Пусть \(k\) --- размер candidate set, а \(\rho\) --- число раундов итерационного локального
поиска.

\begin{enumerate}[label=\arabic*.]
  \item Предварительное построение candidate sets требует для каждой вершины вычислить
  расстояния до остальных вершин и частично отсортировать их. Консервативная оценка:
  \(O(n^2 \log k)\).
  \item Один проход candidate-guided 2-opt по маршруту длины \(q_k+2\) имеет стоимость
  \(O(q_k k)\), после чего выполняется cleanup-проход полного 2-opt со стоимостью
  \(O(q_k^2)\).
  \item На \(\rho\) раундах ILS суммарная стоимость внутримаршрутного улучшения составляет
  \[
  O\!\left(\rho \sum_{k=1}^{m} (q_k k + q_k^2)\right).
  \]
  \item Межмаршрутный этап в текущей реализации пересчитывает целевую функцию всего
  набора маршрутов при каждой проверке обмена. Поэтому одна полная внешняя итерация
  этого этапа в худшем случае имеет кубическую оценку \(O(n^3)\).
\end{enumerate}

Последний пункт важен для интерпретации результатов: \texttt{lkh-wrapper-v2} существенно
улучшает качество решения, но его межмаршрутная фаза все еще остается вычислительно
дорогой. Следовательно, дальнейшее развитие алгоритма логично связывать не только с новыми
эвристическими идеями, но и с оптимизацией стоимости оценки межмаршрутных обменов.

\section{Почему \texttt{lkh-wrapper-v2} должен превосходить baseline}

С точки зрения структуры поиска преимущество \texttt{lkh-wrapper-v2} над базовыми методами
опирается на три обстоятельства.

Во-первых, алгоритм использует более информативную конструктивную фазу. В
\texttt{rand+nn} распределение клиентов между маршрутами фактически случайно, а в
\texttt{2opt+greed} оно опирается только на текущий локальный шаг. В \texttt{lkh-wrapper-v2}
выбор следующего клиента зависит не только от расстояния до него, но и от оценки
дальнейшего продолжения маршрута и удаленности от депо.

Во-вторых, локальный поиск ведется в более сильной окрестности. Полный 2-opt уже сам по себе
сильнее простого ближайшего соседа \cite{croes1958}, а candidate-guided схема и крупные
возмущения позволяют быстрее переходить между качественными локальными минимумами, что
согласуется с идеями Lin--Kernighan и Helsgaun \cite{lin1973,helsgaun2000}.

В-третьих, \texttt{lkh-wrapper-v2} сочетает внутримаршрутные и межмаршрутные улучшения.
Это особенно важно для mTSP, где качество решения определяется не только порядком посещения
вершин внутри маршрута, но и самим распределением клиентов между агентами
\cite{bektas2006,russell1977}. Поэтому ограничиваться только 2-opt внутри уже зафиксированных
маршрутов недостаточно.

\section{Вывод}

Проведенное обоснование позволяет сделать три вывода.

\begin{enumerate}[label=\arabic*.]
  \item Все реализованные в проекте алгоритмы строят допустимые решения задачи mTSP.
  \item Алгоритмы, использующие локальные улучшения с приемом только строго улучшающих
  шагов, завершаются за конечное число итераций.
  \item Улучшения, реализованные в \texttt{grasp} и особенно в \texttt{lkh-wrapper-v2},
  имеют не только эмпирический, но и структурный математический смысл: они расширяют
  класс рассматриваемых решений и усиливают локальный поиск без нарушения допустимости.
\end{enumerate}

Следовательно, переход от baseline-эвристик к \texttt{lkh-wrapper-v2} можно считать
обоснованным не только по вычислительным экспериментам, но и по внутренней логике
алгоритма: улучшение обеспечивается более сильной конструктивной фазой, более содержательной
локальной окрестностью и явной межмаршрутной дооптимизацией.

\begin{thebibliography}{99}

\bibitem{bektas2006}
Bektas T. The multiple traveling salesman problem: an overview of formulations and solution procedures //
Omega. 2006. Vol.~34, No.~3. P.~209--219. DOI: \href{https://doi.org/10.1016/j.omega.2004.10.004}{10.1016/j.omega.2004.10.004}.

\bibitem{croes1958}
Croes G.~A. A Method for Solving Traveling-Salesman Problems //
Operations Research. 1958. Vol.~6, No.~6. P.~791--812. DOI:
\href{https://doi.org/10.1287/opre.6.6.791}{10.1287/opre.6.6.791}.

\bibitem{lin1973}
Lin S., Kernighan B.~W. An Effective Heuristic Algorithm for the Traveling-Salesman Problem //
Operations Research. 1973. Vol.~21, No.~2. P.~498--516. DOI:
\href{https://doi.org/10.1287/opre.21.2.498}{10.1287/opre.21.2.498}.

\bibitem{helsgaun2000}
Helsgaun K. An Effective Implementation of the Lin--Kernighan Traveling Salesman Heuristic //
European Journal of Operational Research. 2000. Vol.~126, No.~1. P.~106--130. DOI:
\href{https://doi.org/10.1016/S0377-2217(99)00284-2}{10.1016/S0377-2217(99)00284-2}.

\bibitem{feo1995}
Feo T.~A., Resende M.~G.~C. Greedy Randomized Adaptive Search Procedures //
Journal of Global Optimization. 1995. Vol.~6, No.~2. P.~109--133. DOI:
\href{https://doi.org/10.1007/BF01096763}{10.1007/BF01096763}.

\bibitem{russell1977}
Russell R.~A. Technical Note---An Effective Heuristic for the M-Tour Traveling Salesman Problem with Some Side Conditions //
Operations Research. 1977. Vol.~25, No.~3. P.~517--524. DOI:
\href{https://doi.org/10.1287/opre.25.3.517}{10.1287/opre.25.3.517}.

\end{thebibliography}

\end{document}
